\documentclass[11pt]{article}
\usepackage[a4paper,margin=2.4cm]{geometry}
\usepackage{booktabs,graphicx,amsmath,natbib,xcolor}
\usepackage[hidelinks]{hyperref}
\renewcommand{\thetable}{S\arabic{table}}
\renewcommand{\thefigure}{S\arabic{figure}}
\renewcommand{\thesection}{S\arabic{section}}
\title{Supplementary Information\\[3pt]\large Non-linguistic communication of abstract rules emerges from reward}
\author{Cristiano Capone\\ \small Quantitative Health Models and Bioinspired Technologies Lab, Istituto Superiore di Sanit\`a, Rome, Italy}
\date{}
\begin{document}
\maketitle

\section{Which held-out programs are new}\label{s:novelty}
A held-out program tests composition only if the transformation it computes is absent from what training presents.
Three routes make it present. Algebra: the geometric operations compose into one another (a vertical reflection
followed by a horizontal one \emph{is} the rotation by 180 degrees, on every grid), and a colour operation commutes
with a spatial one, so a held-out pair can equal a training single or a trained pair in the other order. Training
triples: a held-out pair can equal a trained triple whose last operation is redundant (drawing a border of colour 1
and then applying gravity to the left equals the trained triple that adds a crop). Augmentation: every training
episode is transformed by one of the eight dihedral transforms $T$ and a permutation $s$ of the colours, applied to
demonstrations and query alike, so a training program $P$ is trained as $sTPT^{-1}s^{-1}$. A held-out program equal
to such a conjugate has been trained although it appears in no list of training programs: shifting up and then
applying gravity downwards, for instance, is a rotated copy of the trained pair that shifts right and applies
gravity to the left, and drawing a border of colour 7 and then swapping colours 4 and 8 is a recoloured copy of
the trained pair that swaps colours 9 and 7 and then draws a border of colour 1 (the two steps commute).

We compared every held-out program with every training program (single operations, trained pairs, trained
self-compositions and trained triples) under all eight dihedral transforms and every colour permutation (script
\texttt{audit\_group.py}). The colours an operation names were found by testing its commutation with all colour
transpositions. The colours a program names are not always the colours that matter: in \emph{gravity left}, then
\emph{recolour 9 to 3}, then \emph{recolour 9 to 1}, the last step is idle, because the first leaves no 9, so this
held-out triple computes a trained pair even though it names one colour more. Colour permutations are therefore
enumerated as the injective maps from the colours the training program names into the colours either program
names, all of them, the identity first, and nothing is rejected because the two programs name different colours. Equality was
tested on 19 grids built to separate operations (all cells distinct; all nine colours; non-square shapes) and on
2\,000 fresh random grids. Operations that break ties by colour index or by position (majority fill; in library B
keeping the most or least common colour, ranking colours, marking the densest row or column, sorting rows or
columns) make a conjugate differ from the original only where such a step meets a tie; grids on which either
program meets a tie were left out, and a match required at least 300 tie-free grids, so equality here is equality
up to tie-breaking. Candidates are screened only by output shape and by which cells are empty, neither of which a
colour permutation changes, and only on grids without a tie. Matches are verified on these batteries, not proved. Table~\ref{tab:novelty} gives, for every
library, the held-out programs that no training program computes directly, nor through a dihedral transform, nor
through the full augmentation group; the last are called new.

\begin{table}[h]\centering\small\resizebox{\textwidth}{!}{\begin{tabular}{lcccccccc}\toprule
 & \multicolumn{4}{c}{held-out pairs} & \multicolumn{4}{c}{held-out triples} \\ \cmidrule(lr){2-5}\cmidrule(lr){6-9}
library & programs & not direct & not dihedral & new & programs & not direct & not dihedral & new \\ \midrule
A (30 operations) & 212 & 91 & 32 & 6 & 300 & 231 & 175 & 100 \\
A, first 24 & 134 & 36 & 9 & 6 & 300 & 199 & 107 & 80 \\
A, first 17 & 66 & 15 & 0 & 0 & 300 & 128 & 49 & 49 \\
A, random 17 & 67 & 14 & 10 & 6 & 300 & 167 & 127 & 85 \\
A without 7 confusable & 124 & 41 & 14 & 5 & 300 & 207 & 134 & 80 \\
B (17 operations) & 68 & 31 & 25 & 15 & 300 & 185 & 177 & 169 \\
B plus 6 near-duplicates & 126 & 66 & 48 & 17 & 300 & 205 & 200 & 178 \\
\bottomrule
\end{tabular}}
\caption{Held-out programs not computed by any training program directly (not direct), nor through a dihedral
transform (not dihedral), nor through any dihedral transform combined with a colour permutation (new), per
library and split. Training programs include trained triples.}\label{tab:novelty}
\end{table}

The audit changes what the held-out pairs can test. In the main library 6 of the 212 held-out pairs are new,
and every one of them occurs as consecutive operations of a training triple; in the libraries built from it, 0 to
6. Held-out triples keep a substantial set of new programs in every library: 100 of 300 in the main library, 49 to 85 in its
subsets, and 169 to 178 in library B, with or without near-duplicate operations. The manuscript therefore tests generalisation on new triples, evaluated
with ten episodes per triple, and reports held-out pairs as combinations absent from the list of training programs
but, with those six exceptions, covered by augmentation. Every checkpoint was evaluated on its library's new
triples and new pairs (twenty episodes per pair; Table~\ref{tab:novelresults}). The contrasts the manuscript
rests on are all preserved on new triples (Section~S8). On the six new pairs of the main library accuracy is
low for every channel, too few programs to carry a conclusion. The
held-out self-compositions (an operation applied twice) are not reported: three of their five programs equal a
training single.

\begin{table}[h]\centering\small\resizebox{\textwidth}{!}{\begin{tabular}{lccccc}\toprule
arm & pairs, all & pairs, novel & triples, all & triples, novel & runs \\ \midrule
A, staged, reward & 31.6 $\pm$ 0.7 & 4.8 $\pm$ 0.7 & 17.2 $\pm$ 0.6 & 15.8 $\pm$ 0.4 & 8 \\
A, joint, reward & 1.5 $\pm$ 0.3 & 1.2 $\pm$ 0.4 & 1.7 $\pm$ 0.1 & 2.3 $\pm$ 0.1 & 6 \\
A, reversed, reward & 3.4 $\pm$ 1.6 & 0.8 $\pm$ 0.0 & 2.2 $\pm$ 0.3 & 3.2 $\pm$ 0.8 & 3 \\
A, joint + information term & 29.1 $\pm$ 0.8 & 6.5 $\pm$ 0.3 & 16.6 $\pm$ 0.4 & 15.0 $\pm$ 0.3 & 8 \\
A, joint, changed-cell reward & 2.2 $\pm$ 0.5 & 2.9 $\pm$ 0.4 & 1.7 $\pm$ 0.2 & 1.9 $\pm$ 0.2 & 2 \\
A, staged, changed-cell reward & 3.3 $\pm$ 1.5 & 2.9 $\pm$ 0.4 & 3.2 $\pm$ 1.2 & 3.2 $\pm$ 1.1 & 2 \\
A, staged, gradient & 39.2 $\pm$ 0.2 & 6.5 $\pm$ 0.5 & 22.3 $\pm$ 0.3 & 19.1 $\pm$ 0.3 & 4 \\
A, joint, gradient & 38.7 $\pm$ 0.4 & 7.2 $\pm$ 0.3 & 21.4 $\pm$ 0.5 & 19.3 $\pm$ 0.6 & 3 \\
A, reversed, gradient & 37.4 $\pm$ 0.3 & 8.3 $\pm$ 0.8 & 21.2 $\pm$ 0.2 & 18.4 $\pm$ 0.2 & 2 \\
A, continuous & 39.1 $\pm$ 0.7 & 8.3 $\pm$ 0.3 & 21.4 $\pm$ 0.5 & 19.1 $\pm$ 0.4 & 4 \\
A, reward, solver perception & 21.0 $\pm$ 2.7 & 5.7 $\pm$ 0.5 & 11.4 $\pm$ 1.5 & 11.0 $\pm$ 0.9 & 6 \\
A, gradient, solver perception & 37.9 $\pm$ 0.6 & 6.7 $\pm$ 0.3 & 22.0 $\pm$ 0.2 & 19.4 $\pm$ 0.3 & 4 \\
A, reward, end to end & 0.6 $\pm$ 0.0 & 0.0 $\pm$ 0.0 & 0.8 $\pm$ 0.2 & 1.7 $\pm$ 0.1 & 2 \\
A, gradient, end to end & 14.1 $\pm$ 1.6 & 7.1 $\pm$ 0.4 & 8.5 $\pm$ 1.3 & 7.7 $\pm$ 1.8 & 2 \\
A24, staged / joint, reward & 18.9 $\pm$ 10.8 / 5.1 $\pm$ 0.4 & 6.2 $\pm$ 4.6 / 2.1 $\pm$ 2.1 & 8.2 $\pm$ 4.7 / 2.3 $\pm$ 0.3 & 8.2 $\pm$ 2.9 / 5.7 $\pm$ 0.8 & 2 / 2 \\
A17, staged / joint, reward & 15.2 $\pm$ 8.2 / 0.9 $\pm$ 0.3 & -- / -- & 10.8 $\pm$ 5.6 / 2.4 $\pm$ 1.1 & 7.9 $\pm$ 3.8 / 2.3 $\pm$ 0.1 & 2 / 2 \\
A17r, staged / joint, reward & 27.9 $\pm$ 3.4 / 21.5 $\pm$ 9.0 & 15.4 $\pm$ 2.9 / 15.4 $\pm$ 2.1 & 14.3 $\pm$ 2.8 / 13.7 $\pm$ 3.2 & 10.6 $\pm$ 0.4 / 8.9 $\pm$ 0.6 & 2 / 2 \\
A$-$conf, staged / joint, reward & 36.0 $\pm$ 2.0 / 3.1 $\pm$ 0.1 & 16.5 $\pm$ 0.5 / 2.0 $\pm$ 0.0 & 20.3 $\pm$ 1.2 / 3.5 $\pm$ 0.2 & 8.7 $\pm$ 0.6 / 3.1 $\pm$ 0.3 & 2 / 2 \\
B, staged / joint, reward & 12.1 $\pm$ 0.3 / 12.8 $\pm$ 0.1 & 8.5 $\pm$ 0.2 / 8.8 $\pm$ 0.2 & 13.2 $\pm$ 0.1 / 14.0 $\pm$ 0.2 & 6.7 $\pm$ 0.1 / 7.2 $\pm$ 0.1 & 2 / 2 \\
B, joint + information term & 12.8 $\pm$ 0.1 & 8.3 $\pm$ 0.0 & 13.8 $\pm$ 0.3 & 7.3 $\pm$ 0.0 & 2 \\
B+dup, staged / joint, reward & 9.3 $\pm$ 0.2 / 10.2 $\pm$ 0.2 & 3.7 $\pm$ 0.1 / 3.8 $\pm$ 0.3 & 10.9 $\pm$ 0.1 / 11.1 $\pm$ 0.1 & 7.3 $\pm$ 0.1 / 7.6 $\pm$ 0.1 & 2 / 2 \\
\bottomrule
\end{tabular}}
\caption{Exact match (\%) on all held-out programs and on the new ones, mean $\pm$ s.e.m. over runs (a single
value where there is one run). For the smaller libraries, staged / joint.}\label{tab:novelresults}
\end{table}

\section{Decoding the rule from the message}\label{s:decoding}
A common summary of how much a code says about its referents is the normalised mutual information between rule
family and message, estimated from the joint distribution in the sample. That estimate is 1 whenever messages are
nearly all distinct, because each message then identifies its rule \emph{in the sample} whatever the real
relationship: independent random messages assigned to 239 families $\times$ 10 episodes reach it. We therefore
decode on held-out episodes instead. The episodes of each family are split in two halves; a decoder
fitted on one half is scored on the other, and the halves are swapped. Exact lookup succeeds only if the
same message recurs for the family; nearest neighbour in Hamming distance succeeds if the family's messages
occupy a consistent region of message space. On independent random messages nearest-neighbour decoding
gives 0.003 against a chance level of 0.004.

\begin{table}[h]\centering\small\begin{tabular}{rcccccc}\toprule
 & \multicolumn{2}{c}{gradient} & \multicolumn{2}{c}{reward} & & \\ \cmidrule(lr){2-3}\cmidrule(lr){4-5}
families & nearest & lookup & nearest & lookup & chance & null, 95th pct. \\ \midrule
30 & 0.725 $\pm$ 0.032 & 0.602 $\pm$ 0.021 & 0.308 $\pm$ 0.046 & 0.317 $\pm$ 0.035 & 0.033 & 0.059 \\
60 & 0.630 & 0.397 & 0.187 $\pm$ 0.028 & 0.170 $\pm$ 0.020 & 0.017 & 0.027 \\
120 & 0.583 & 0.303 & 0.136 $\pm$ 0.035 & 0.144 $\pm$ 0.035 & 0.008 & 0.015 \\
239 & 0.424 $\pm$ 0.002 & 0.219 $\pm$ 0.010 & 0.091 $\pm$ 0.000 & 0.090 $\pm$ 0.005 & 0.004 & 0.007 \\
\bottomrule
\end{tabular}
\caption{Decoding of the rule family from the message on held-out episodes, ten episodes per family; the
null is the 95th percentile over 200 permutations of the family labels in each fitted half, scored with the same
two-fold average as the data, the larger of
the two channels.}\label{tab:decoding}
\end{table}

The gradient code decodes four to five times better than the reward code at every size, and far above the
null. For the reward code exact lookup and nearest neighbour agree: it uses fixed labels. For the gradient
code nearest neighbour is roughly twice lookup at 60 families and beyond: the same family rarely receives the
identical message but its messages cluster, which is why its purity falls with the number of families while
its decodability stays high.

\begin{table}[h]\centering\small\begin{tabular}{lccc}\toprule
239 families (unless stated) & nearest & lookup & runs \\ \midrule
reward, baseline & 0.091 $\pm$ 0.000 & 0.090 $\pm$ 0.005 & 2 \\
message length 16 & 0.093 $\pm$ 0.005 & 0.069 $\pm$ 0.011 & 2 \\
vocabulary 128 & 0.060 $\pm$ 0.021 & 0.061 $\pm$ 0.016 & 2 \\
information term 0.1 & 0.090 $\pm$ 0.005 & 0.115 $\pm$ 0.004 & 2 \\
information term 0.3 & 0.136 & 0.156 & 1 \\
reward, 319 families, term 0.1 & 0.067 & 0.085 & 1 \\
gradient reference & 0.424 $\pm$ 0.002 & 0.219 $\pm$ 0.010 & 2 \\
\bottomrule
\end{tabular}
\caption{What changes the decodability of a reward-trained code at 239 families.}\label{tab:decomposition}
\end{table}

Longer messages add distinct messages and no decodable information, and a larger vocabulary lowers it.
The information term does not change it at weight 0.1 and raises it modestly at 0.3, from one seed. The
in-sample mutual information would have suggested that the term relieves naming saturation at weight 0.1;
decoding on held-out episodes shows that it does not.

\section{No evidence of serial execution}
The executor applies the eight symbols one after another, so its state after the first symbols of a
two-operation rule could come to resemble the intermediate grid. We decoded the executor's state after every
symbol on held-out pairs and took the best-matching state against the true intermediate grid (union cell
accuracy): 0.695 $\pm$ 0.009 for the continuous channel, 0.703 $\pm$ 0.006 for gradient and 0.737 $\pm$ 0.010
for reward (4, 4 and 5 runs). Copying the input matches the intermediate grid at 0.670 and the final output at
0.647. The best state therefore exceeds both references by 3--9 points; we did not test whether the maximum over
seven states of a control executor would give the same margin, and no state approaches the intermediate grid. We read this as no
evidence of serial execution, not as evidence against it.

\section{Message swaps}
Applying to a query of $A\!\circ\!B$ the message the sender produced for $A\!\circ\!C$ (a donor sharing
the first operation) gives the executor an output that matches the true intermediate grid of
$A\!\circ\!B$ at 0.70 cell accuracy in both discrete channels, against 0.65 with the query's own
message and 0.67 for the input itself. The donor message therefore carries the shared first operation
to a small but consistent degree, without the executor's state ever representing it as a separate
step (Section~S3).


\section{Name stability, and the accounts the measurements did not support}\label{s:falsified}
Six accounts of the curriculum effect, or of which libraries collapse, were tested; none was supported by the
measurements, each within the range of the manipulation we could run. They are listed so that they are not
re-run. The section ends with the measure that does account for competence.
(1)~Credit assignment: giving each symbol its own credit from the gain in the executor's state leaves joint
training where it was.
(2)~The number of operations: two sets of seventeen behave oppositely.
(3)~Perceptual confusability: removing the seven operations the frozen perception confuses leaves library A
collapsing, adding six near-duplicates to library B does not collapse it, and in the self-supervised
perception library A has only three confusable operations and collapses completely.
(4)~The value of the message to the executor, and the partial credit one operation earns when answered with
another's output: neither orders the libraries.
(5)~A race between the sender's policy sharpening and the executor's learning: the sender's entropy does
crash within 800 steps in the collapsing libraries and not in the others, but keeping it exploring longer adds
little information to the code and little accuracy (Table~\ref{tab:infocomp}), and forcing library B's sender to
sharpen early does not collapse it.
(6)~Name stability (Table~\ref{tab:stability}): the hypothesis that a useful code needs names that stay fixed
long enough for the executor to learn them. We logged the modal message of every single operation every 500
steps. The most competent code, from staged experience, has the highest turnover of any arm, at a purity
comparable to the changed-cell code that fails, whose turnover lies between the two competent arms. A pressure
towards the sender's own recent average (the Kullback--Leibler divergence to an exponential moving average of
its parameters) did not stabilise names: at weight 0.1 it changed nothing, at weight 1.0 it softened the policy
and raised the measured turnover. Applied to plain joint reward at weight 1.0, the same pressure lifts accuracy on
held-out pairs from 0.000--0.014 to 0.12--0.15 by slowing the sender's early commitment; it is a third, weaker
remedy for collapse, not a stabiliser of names.

\begin{table}[h]\centering\small\begin{tabular}{lccc}\toprule
arm (library A, two seeds) & name turnover after 16\,000 steps & purity & exact, held-out pairs \\ \midrule
staged & 0.31 / 0.24 & 0.81 / 0.78 & 0.307 / 0.302 \\
joint + information term & 0.06 / 0.09 & 0.91 / 0.87 & 0.274 / 0.292 \\
joint, changed-cell reward & 0.13 / 0.15 & 0.83 / 0.82 & 0.024 / 0.024 \\
\quad + stability pressure 0.1 & 0.07 / 0.13 & 0.86 / 0.80 & 0.047 / 0.028 \\
\quad + stability pressure 1.0 & 0.24 / 0.25 & 0.60 / 0.60 & 0.080 / 0.080 \\
joint, plain reward & 0.05 / 0.02 & 0.97 / 0.98 & 0.000 / 0.014 \\
\quad + stability pressure 1.0 & 0.29 / 0.27 & 0.56 / 0.58 & 0.123 / 0.146 \\
\bottomrule
\end{tabular}
\caption{Name turnover (fraction of the thirty single operations whose modal message changed between
evaluations 500 steps apart, averaged after step 16\,000), purity and accuracy, two seeds per arm.}\label{tab:stability}
\end{table}

\paragraph{Information in the code and competence.} Counted as distinct messages per operation, a reward paying only
for changed cells gives the most differentiated code of any joint run (1.07--1.13 messages per operation) while
solving almost nothing, which would suggest that a differentiated code can be useless. The count is misleading: it
includes variation of the message within an operation. Counted as names (distinct modal messages) the changed-cell
code has 13--16 for thirty operations, one of them shared by 15--17 operations, and its messages identify a new
triple in 10--11\% of episodes, against 8--10\% for the collapsed code and 24--35\% for the staged one.
Table~\ref{tab:infocomp} gives, for every condition on library A, accuracy on the new triples, how often the most
probable message identifies the triple and the operation on held-out episodes (nearest-neighbour decoding; chance
1.0\% and 3.3\%), and the number of names. Over the 50 reward-trained runs accuracy rises with triple identification
(Spearman $\rho = 0.92$, 95\% bootstrap interval 0.84--0.96) and with operation identification ($\rho = 0.80$,
0.67--0.88); including the six runs trained with a pressure towards the sender's recent policy, $\rho = 0.89$
(0.80--0.94) and 0.73 (0.62--0.81). That pressure keeps the policy soft (purity 0.53--0.60 at weight 1.0), so the
executor is trained on noisy samples, and the most probable message, which is what the decoder reads, overstates
what the executor saw. Within one condition the relation rests on few runs: over the six staged runs with the solver's
perception $\rho = 0.83$, and inside the four libraries built from library A the run whose message identifies
triples more often is the more accurate in 19 of 24 pairs of runs. The relation is a correlation across
manipulations, not a demonstration that missing information is the only cause of failure.

\begin{table}[h]\centering\small\resizebox{\textwidth}{!}{\begin{tabular}{lccccc}\toprule
condition (library A) & runs & exact, new triples (\%) & triple identified (\%) & operation identified (\%) & names (of 30) \\ \midrule
\multicolumn{6}{l}{\emph{reward channel}} \\
\quad staged & 8 & 13.1--17.1 & 23.9--35.1 & 53.3--71.8 & 17--25 \\
\quad joint & 8 & 2.1--2.6 & 8.2--10.3 & 18.0--23.5 & 6--8 \\
\quad reversed & 3 & 1.9--4.7 & 9.3--14.8 & 21.2--23.3 & 7--8 \\
\quad joint + information term & 8 & 13.6--16.4 & 23.9--30.1 & 60.8--74.0 & 19--24 \\
\quad staged + information term & 1 & 12.2 & 30.2 & 64.0 & 20 \\
\quad joint, changed-cell reward & 2 & 1.8--2.1 & 10.1--10.7 & 39.7--40.3 & 13--16 \\
\quad staged, changed-cell reward & 2 & 2.1--4.4 & 9.3--12.1 & 41.7--47.3 & 14--15 \\
\quad joint, slower exploration & 2 & 4.3--4.6 & 16.1--17.6 & 27.0--28.5 & 11--12 \\
\quad staged, per-symbol credit & 3 & 4.5--15.5 & 9.9--33.8 & 23.5--71.5 & 7--24 \\
\quad joint, per-symbol credit & 3 & 2.2--2.6 & 10.5--11.9 & 18.0--24.3 & 6--8 \\
\quad staged, solver perception & 6 & 9.1--15.2 & 19.7--28.7 & 31.5--63.2 & 11--20 \\
\quad joint, solver perception & 4 & 2.7--4.7 & 12.8--17.0 & 18.8--21.2 & 6--8 \\
\multicolumn{6}{l}{\emph{reward channel, policy held soft by a stability pressure}} \\
\quad joint, changed-cell + stability 0.1 & 2 & 2.3--3.0 & 9.1--10.7 & 39.8--40.0 & 13--14 \\
\quad joint, changed-cell + stability 1.0 & 2 & 6.7--7.0 & 20.9--24.4 & 72.0--73.2 & 26--26 \\
\quad joint + stability 1.0 & 2 & 8.6--9.2 & 36.9--37.1 & 75.2--79.5 & 28--29 \\
\multicolumn{6}{l}{\emph{gradient channel}} \\
\quad gradient, staged & 4 & 18.4--19.5 & 47.1--53.2 & 87.3--89.8 & 28--29 \\
\quad gradient, joint & 3 & 18.1--20.1 & 49.7--55.6 & 85.2--86.3 & 28--29 \\
\quad gradient, reversed & 2 & 18.2--18.5 & 51.8--52.0 & 87.0--87.3 & 28--28 \\
\bottomrule
\end{tabular}
}
\caption{Competence and information in the code, library A, self-supervised perception unless stated: range
over runs of exact match on the new triples, of how often the most probable message identifies the triple (ten
episodes per triple) and the operation (twenty episodes per single operation) on held-out episodes, and of the
number of names for the thirty single operations. Joint experience includes two runs of the
same configuration from the stability experiment.}\label{tab:infocomp}
\end{table}

\paragraph{Co-adaptation.} The objective of a signalling game can fail through its information component, when the
message does not identify the input, or through its co-adaptation component, when the listener has not adapted to
the sender \cite{rita2022}. To separate the two we froze trained senders and trained a new executor for 8\,000
steps on their most probable messages, and a second one on the same messages permuted across the episodes of each
batch (Table~\ref{tab:coadapt}). The original executor was trained for 32\,000 steps, so a new executor that comes
close to it is learning the code quickly. From staged and information-term codes a new executor reaches
11--12\% of new triples, against 12.6--13.6\% for the original executors on the same episodes; from the collapsed,
changed-cell and slower-exploration codes at most 3\%; from permuted messages 0.4--0.5\%. The gradient code's
messages give 15\% and the continuous channel's 16\%. Where the policy was held soft by the stability pressure at
weight 1.0, the new executor exceeds the original, clearly on single operations (31 and 41\% against 16 and 30\%,
150 episodes each) and less on new triples (6.7 and 10.6\% against 5.4 and 7.8\%), as a partial failure of
coupling predicts: the original executor learned from noisy samples of a code whose most probable messages are
informative. Each of these two conditions rests on one sender, and a new executor also gains a couple of points on
some uninformative codes (staged experience with changed-cell reward: 4.2\% against 2.9\%).

\begin{table}[h]\centering\small\resizebox{\textwidth}{!}{\begin{tabular}{lcccccc}\toprule
 & \multicolumn{3}{c}{new triples (\%)} & \multicolumn{3}{c}{single operations (\%)} \\ \cmidrule(lr){2-4}\cmidrule(lr){5-7}
sender (library A) & original & fresh & blind & original & fresh & blind \\ \midrule
reward, staged & 13.6 / 13.6 & 11.5 / 11.0 & 0.4 / -- & 58.0 / 54.0 & 46.7 / 42.7 & 0.0 / -- \\
reward, joint + information term & 12.6 / 13.5 & 11.7 / 11.9 & 0.5 / -- & 50.7 / 50.7 & 44.0 / 44.7 & 0.0 / -- \\
reward, joint & 0.7 / 1.0 & 1.2 / 0.9 & 0.4 / -- & 2.0 / 2.7 & 3.3 / 2.0 & 0.0 / -- \\
reward, joint, changed-cell reward & 1.8 / 1.2 & 2.3 / 3.2 & 0.5 / -- & 5.3 / 4.7 & 8.0 / 10.7 & 0.0 / -- \\
reward, staged, changed-cell reward & 2.9 & 4.2 & -- & 7.3 & 11.3 & -- \\
reward, joint, slower exploration & 2.8 & 2.8 & -- & 9.3 & 10.0 & -- \\
reward, joint, changed-cell + stability 1.0 (soft policy) & 5.4 & 6.7 & 0.5 & 16.0 & 31.3 & 0.0 \\
reward, joint + stability 1.0 (soft policy) & 7.8 & 10.6 & -- & 30.0 & 40.7 & -- \\
gradient, staged & 17.9 / 17.1 & 15.1 / 14.9 & -- / -- & 77.3 / 77.3 & 67.3 / 66.0 & -- / -- \\
continuous, staged & 18.1 & 15.7 & -- & 62.0 & 58.0 & -- \\
\bottomrule
\end{tabular}
}
\caption{Co-adaptation test, library A, self-supervised perception: exact match of the executor trained with the
sender (original), of a new executor trained for 8\,000 steps on the frozen sender's most probable messages
(fresh), and of a new executor trained on those messages permuted across the episodes of each batch (blind), on
the new triples (ten episodes per triple) and on 150 single-operation episodes; one value per sender, separated
by slashes.}\label{tab:coadapt}
\end{table}

\section{Transfer to real ARC-1 rules}
We trained the three channels on the 320 training rule families of RE-ARC (the procedural generators of
the ARC-1 training tasks), at the easy end of each generator's difficulty range, canvas 16, and
evaluated on 80 held-out families. Because 11\% of easy queries have output equal to input, we report
exact match on non-trivial queries (output $\neq$ input), whose identity floor is zero, and count the
held-out families with at least one such solve. Table~\ref{tab:rearc}. The discrete gradient code
carries real rules at 90\% of the continuous channel's level; the reward code at a third. Only three
of the eighty families are solved with an output that is not a single uniform colour. On the ARC-1
evaluation set the exact match is 0.000 for every channel; at the generators' full difficulty the
continuous channel reaches 0.016 on held-out families.

\begin{table}[h]\centering\small
\resizebox{\textwidth}{!}{\begin{tabular}{lccc}\toprule
 & exact, non-trivial queries, held-out families & families with $\geq 1$ solve & distinct modal messages / 319 rules \\ \midrule
continuous ($n=2$) & 0.070 / 0.070 & 9 / 80 & -- \\
gradient, frozen perception ($n=2$) & 0.063 / 0.063 & 10 / 80 & 304 / 297 \\
gradient, from scratch & 0.046 & 7 / 80 & 261 \\
reward, frozen perception ($n=2$) & 0.021 / 0.025 & 5 / 80 & 63 / 64 \\
full-difficulty control (continuous, canvas 30) & 0.016 & -- & -- \\ \bottomrule
\end{tabular}}
\caption{Transfer to held-out ARC-1 rule families (RE-ARC, easy difficulty, 20\,000 steps).}\label{tab:rearc}
\end{table}

\section{Comparison with human descriptions of the same rules}
LARC \citep{acquaviva2022} contains natural-language descriptions of every ARC-1 training task, written
by people who solved it, together with the success of other people who tried to solve a new instance
from the description alone. Because RE-ARC and LARC index the same 400 tasks, two comparisons are
possible.

\textbf{Transmissibility.} For each rule, the size of the cluster of rules sharing its modal message
in our code, against human build success, description length and describer confidence
(Table~\ref{tab:larc1}). The gradient codes name nearly every rule uniquely, so cluster size has no
variance. The reward code merges 341 of 399 rules into shared messages, and the rules it merges are
the ones people transmit \emph{more} successfully, with shorter descriptions and higher confidence: the
opposite of the hypothesis that the code merges what people find hard, and consistent with a generic
message being partly right for simple transformations. The effects are small ($\rho \le 0.19$) and
confounded by difficulty, since LARC's participants saw the original instances and our system the easy
regenerated ones.

\begin{table}[h]\centering\small
\resizebox{\textwidth}{!}{\begin{tabular}{lcccc}\toprule
code & rules uniquely named & $\rho$(cluster, build success) & $\rho$(cluster, description length) & $\rho$(cluster, confidence) \\ \midrule
gradient, frozen perception & 372 / 399 & $-0.004$ ($p=.94$) & $-0.030$ ($p=.55$) & $+0.063$ ($p=.21$) \\
reward, frozen perception & 58 / 399 & $+0.144$ ($p=.004$) & $-0.137$ ($p=.006$) & $+0.189$ ($p<.001$) \\
gradient, from scratch & 299 / 399 & $-0.004$ ($p=.94$) & $-0.027$ ($p=.58$) & $+0.063$ ($p=.21$) \\ \bottomrule
\end{tabular}}
\caption{Rule-level correlation (Spearman) between the size of a rule's message cluster and human
transmissibility measures from LARC ($n = 399$ rules).}\label{tab:larc1}
\end{table}

\textbf{Category alignment.} Whether two rules that share a message in our code are described more
similarly by people than two rules that do not: Mantel correlations between the same-message (or
Hamming-similarity) matrix of the code and the cosine similarity of the pooled descriptions (TF--IDF
over 317 tasks with human builds) or of their concept-tag profiles (LARC's 17 annotated concepts, 76
tasks), with 2\,000 label permutations, with and without residualisation on human build success and
description length. Table~\ref{tab:larc2}. All correlations are below 0.06; the only nominally
significant ones (Hamming similarity against text, $\rho = 0.04$--$0.06$) vanish with the giant
cluster removed or are within the range expected from 317 tasks. The categories the reward code forms
are not the categories people describe.

\begin{table}[h]\centering\small
\resizebox{\textwidth}{!}{\begin{tabular}{llcccc}\toprule
code & human measure & same message & Hamming similarity & same message, giant cluster removed & $n$ \\ \midrule
reward s0 & text & $+0.008$ & $+0.040$ ($p=.015$) & $-0.002$ & 317 \\
reward s1 & text & $+0.011$ & $+0.036$ ($p=.037$) & $-0.015$ & 317 \\
reward s0 & concept tags & $-0.023$ & $-0.001$ & $-0.108$ & 76 \\
reward s1 & concept tags & $+0.022$ & $+0.034$ & $-0.063$ & 76 \\
gradient s0 & text & $+0.014$ & $+0.056$ ($p<.001$) & $+0.008$ & 317 \\
gradient s1 & text & $+0.014$ & $+0.049$ ($p<.001$) & $+0.010$ & 317 \\
gradient s0 & concept tags & $+0.007$ & $+0.012$ & $+0.008$ & 76 \\
gradient s1 & concept tags & $+0.013$ & $+0.061$ ($p=.041$) & $+0.012$ & 75 \\ \bottomrule
\end{tabular}}
\caption{Mantel correlations between the code's rule-similarity matrices and human description
similarity. Partial correlations controlling for difficulty differ by $< 0.01$.}\label{tab:larc2}
\end{table}


\bibliographystyle{plainnat}
\bibliography{refs}
\section{Statistics for every contrast}\label{s:stats}
Every contrast between conditions reported in the paper, with the number of runs in each condition. Accuracy
contrasts use exact match on each library's new triples, ten episodes per triple. The first three contrasts are
the primary ones; the others are exploratory, and a difference whose interval includes zero is not evidence of
equality. The difference
and its 95\% confidence interval come from a hierarchical bootstrap that resamples runs within each condition
and, jointly for the conditions compared, the held-out programs (10\,000 resamples), with the two-sided bootstrap
$p$-value. Welch's $t$-test and an exact permutation test on run means are given alongside, the latter with the
smallest $p$-value its number of runs allows. Where a condition has a single run, the interval reflects
variation over programs only. Decoding contrasts resample runs and rule families; the human contrast resamples
nothing and uses chains as the unit. Differences in accuracy and decoding are in percentage points.

\begin{table}[h]\centering\small\resizebox{\textwidth}{!}{\begin{tabular}{lcccccc}\toprule
contrast & runs & means & difference [95\% CI] & $p_\mathrm{boot}$ & $p_\mathrm{Welch}$ & $p_\mathrm{perm}$ (min) \\ \midrule
\multicolumn{7}{l}{\emph{Primary contrasts}} \\
\quad reward: staged $-$ joint & 8 / 6 & 15.8 / 2.3 & 13.5 [9.1, 18.4] & $<$0.001 & $<$0.001 & $<$0.001 ($<$0.001) \\
\quad reward: joint + information term $-$ joint & 8 / 6 & 15.0 / 2.3 & 12.7 [8.4, 17.5] & $<$0.001 & $<$0.001 & $<$0.001 ($<$0.001) \\
\quad interaction: (staged $-$ joint), reward $-$ gradient & 8 / 6 / 4 / 3 & -- & 13.8 [9.0, 18.7] & $<$0.001 & -- & -- \\
\multicolumn{7}{l}{\emph{Exploratory: discreteness and perception (staged experience)}} \\
\quad gradient $-$ continuous & 4 / 4 & 19.1 / 19.2 & $-$0.1 [$-$1.5, 1.4] & 0.89 & 0.83 & 0.86 (0.03) \\
\quad reward $-$ continuous & 8 / 4 & 15.8 / 19.2 & $-$3.4 [$-$5.6, $-$1.4] & $<$0.001 & $<$0.001 & 0.002 (0.002) \\
\quad reward $-$ gradient & 8 / 4 & 15.8 / 19.1 & $-$3.3 [$-$5.5, $-$1.3] & 0.001 & $<$0.001 & 0.002 (0.002) \\
\quad reward: self-supervised $-$ end-to-end perception & 8 / 2 & 15.8 / 1.7 & 14.1 [9.7, 19.1] & $<$0.001 & $<$0.001 & 0.02 (0.02) \\
\quad gradient: self-supervised $-$ end-to-end perception & 4 / 2 & 19.1 / 7.7 & 11.3 [6.8, 16.6] & $<$0.001 & 0.09 & 0.07 (0.07) \\
\quad reward: self-supervised $-$ solver perception & 8 / 6 & 15.8 / 11.0 & 4.8 [1.7, 8.5] & 0.001 & 0.002 & $<$0.001 ($<$0.001) \\
\multicolumn{7}{l}{\emph{Exploratory: order of experience}} \\
\quad reward: staged $-$ reversed & 8 / 3 & 15.8 / 3.2 & 12.6 [8.0, 17.4] & $<$0.001 & $<$0.001 & 0.006 (0.006) \\
\quad gradient: staged $-$ joint & 4 / 3 & 19.1 / 19.3 & $-$0.3 [$-$1.8, 1.4] & 0.70 & 0.71 & 0.66 (0.03) \\
\quad gradient: staged $-$ reversed & 4 / 2 & 19.1 / 18.4 & 0.7 [$-$0.6, 2.2] & 0.33 & 0.08 & 0.20 (0.07) \\
\multicolumn{7}{l}{\emph{Exploratory: information term and changed-cell reward (reward channel)}} \\
\quad joint + information term (weight 0.1 only) $-$ joint & 6 / 6 & 14.8 / 2.3 & 12.5 [8.2, 17.4] & $<$0.001 & $<$0.001 & 0.002 (0.002) \\
\quad joint + information term $-$ staged & 8 / 8 & 15.0 / 15.8 & $-$0.8 [$-$2.9, 1.2] & 0.44 & 0.14 & 0.15 ($<$0.001) \\
\quad joint, changed-cell reward $-$ joint & 2 / 6 & 1.9 / 2.3 & $-$0.3 [$-$2.0, 1.4] & 0.66 & 0.20 & 0.14 (0.04) \\
\quad joint + information term $-$ joint, changed-cell reward & 8 / 2 & 15.0 / 1.9 & 13.0 [9.0, 17.3] & $<$0.001 & $<$0.001 & 0.02 (0.02) \\
\quad staged, changed-cell reward $-$ staged & 2 / 8 & 3.2 / 15.8 & $-$12.5 [$-$17.3, $-$8.1] & $<$0.001 & 0.03 & 0.02 (0.02) \\
\multicolumn{7}{l}{\emph{Exploratory: other libraries (reward channel)}} \\
\quad A, random 17: staged $-$ joint & 2 / 2 & 10.6 / 8.9 & 1.6 [$-$0.1, 3.8] & 0.06 & 0.17 & 0.33 (0.33) \\
\quad A without confusable operations: staged $-$ joint & 2 / 2 & 8.7 / 3.1 & 5.6 [2.4, 9.3] & $<$0.001 & 0.03 & 0.33 (0.33) \\
\quad A, first 24: staged $-$ joint & 2 / 2 & 8.2 / 5.7 & 2.6 [$-$1.4, 8.1] & 0.33 & 0.53 & 0.67 (0.33) \\
\quad A, first 17: staged $-$ joint & 2 / 2 & 7.9 / 2.3 & 5.5 [$-$0.2, 12.9] & 0.07 & 0.38 & 0.33 (0.33) \\
\quad B: staged $-$ joint & 2 / 2 & 6.7 / 7.2 & $-$0.5 [$-$1.2, 0.0] & 0.08 & 0.09 & 0.33 (0.33) \\
\quad B: joint + information term $-$ joint & 2 / 2 & 7.3 / 7.2 & 0.1 [$-$0.5, 0.6] & 0.86 & 0.50 & 1.00 (0.33) \\
\quad B with near-duplicates: staged $-$ joint & 2 / 2 & 7.3 / 7.6 & $-$0.3 [$-$0.8, 0.0] & 0.07 & 0.24 & 0.33 (0.33) \\
\multicolumn{7}{l}{\emph{Exploratory: collapse diagnostic, distinct messages per held-out rule (run level)}} \\
\quad staged $-$ joint & 8 / 6 & 0.33 / 0.09 & 0.25 [0.20, 0.29] & -- & $<$0.001 & $<$0.001 ($<$0.001) \\
\quad reversed $-$ joint & 3 / 6 & 0.12 / 0.09 & 0.03 [$-$0.04, 0.10] & -- & 0.24 & 0.10 (0.01) \\
\quad joint + information term $-$ joint & 7 / 6 & 0.29 / 0.09 & 0.20 [0.16, 0.25] & -- & $<$0.001 & $<$0.001 ($<$0.001) \\
\quad joint, changed-cell reward $-$ joint & 2 / 6 & 0.11 / 0.09 & 0.02 [$-$0.01, 0.05] & -- & 0.09 & 0.14 (0.04) \\
\multicolumn{7}{l}{\emph{Exploratory: decoding the rule family from the message, 239 families}} \\
\quad gradient $-$ reward, nearest neighbour & 2 / 2 & 42.5 / 9.1 & 33.4 [30.2, 36.6] & $<$0.001 & 0.002 & 0.33 (0.33) \\
\quad gradient: nearest neighbour $-$ exact lookup & 2 / 2 & 42.5 / 21.9 & 20.5 [17.6, 23.4] & $<$0.001 & 0.02 & -- \\
\quad reward: nearest neighbour $-$ exact lookup & 2 / 2 & 9.1 / 8.9 & 0.1 [$-$1.9, 2.0] & 0.85 & 0.82 & -- \\
\quad reward: length 16 $-$ baseline, nearest & 2 / 2 & 9.3 / 9.1 & 0.2 [$-$1.8, 2.3] & 0.82 & 0.70 & 1.00 (0.33) \\
\quad reward: vocabulary 128 $-$ baseline, nearest & 2 / 2 & 6.0 / 9.1 & $-$3.1 [$-$6.5, 0.4] & 0.10 & 0.38 & 0.33 (0.33) \\
\quad reward: information term 0.1 $-$ baseline, nearest & 2 / 2 & 9.0 / 9.1 & $-$0.1 [$-$2.4, 2.1] & 0.90 & 0.84 & 1.00 (0.33) \\
\quad reward: information term 0.3 $-$ baseline, nearest & 1 / 2 & 13.6 / 9.1 & 4.5 [2.0, 7.0] & $<$0.001 & -- & -- \\
\multicolumn{7}{l}{\emph{Exploratory: human iterated learning (chain level)}} \\
\quad final generation, distinct signals per meaning: pressure $-$ none & 4 / 4 & 0.59 / 0.14 & 0.45 [0.17, 0.74] & -- & 0.01 & 0.03 (0.03) \\
\bottomrule
\end{tabular}
}
\caption{Contrasts between conditions. Runs: number of runs in each condition; means in the same order.}\label{tab:stats}
\end{table}

\section{A pretrained model in the same game: a pilot}\label{s:llm}
Everything above concerns a pair of small networks trained from scratch. To see whether the degenerate solution and
its remedy depend on that architecture, we played the same game with a pretrained language model: two copies of
Qwen2.5-0.5B-Instruct, sender and executor, each adapted by low-rank adapters (rank 16 on the attention and
feed-forward projections, 8.8\,M trainable parameters of 503\,M). The sender reads the worked examples as text and
writes eight symbols from a fixed alphabet of 32 tokens that carry no English meaning; the executor reads those
symbols and the query grid and writes the output. Training follows the reward recipe of the main experiments (four
sampled messages per episode, within-group standardised advantage, entropy floor, sender learning rate 0.3 times
the executor's) for 5\,000 steps at batch 16.

Two things differ from the main experiments and limit what the comparison can show. The reward is the accuracy of
the output tokens under teacher forcing, where the main experiments reward the union cell accuracy of the executor's
prediction (both are evaluated by exact match of the whole grid): a generative executor rarely produces a whole
correct grid early in training, and teacher forcing changes both what the reward measures and what it is
conditioned on. And these runs do not augment their episodes, so a held-out program is new to them when no training
program computes it, without rotation, reflection or recolouring: 231 of the 300 held-out triples and 91 of the 212
held-out pairs (Section~\ref{s:novelty}), three episodes each. Accuracy below is exact match on those programs; the
code is described over all 300 held-out triples, from the messages stored during the evaluation.

Table~\ref{tab:llm} reports every run, counted at each stage: started, finished (5\,000 steps), non-finite and
averaged. These runs train in 16-bit floating point, and in four of the eight information-term runs the adapter
parameters became non-finite late in training (steps 3\,850 to 4\,450). The game then restores its last snapshot with
a fresh optimiser; in all four the run degenerated afterwards, its training accuracy falling to zero by step 5\,000.
Each of these runs is therefore evaluated at its last saved snapshot before its first non-finite event (step 4\,000
in three runs, 3\,500 in one; snapshots are saved every 500 steps), a rule fixed before the evaluations were run,
and every other run at its final adapters. What is recovered is a checkpoint, not a stable training run: half of
this arm failed late in training, and we cannot exclude that the instability interacts with the objective.

With joint experience the code collapses in all four runs, exact match tracking the constant-message control almost
exactly in every one (1.3--2.3\% against 1.2--2.3\%, one to nine names for 300 rules, decoding 0.1--0.2\%, at the
permutation null). Staged experience collapses just as completely, in all four of its runs (1.0--1.9\% against
1.0--2.0\%, one to three names, decoding 0.1--0.5\%): unlike the small networks, staged experience does not rescue
this model. Its first stage is 1\,000 steps, too short to separate a real difference from a budget that never learns
the single operations, so this comparison is weaker than the joint-versus-staged contrast of the main experiments;
what four seeds each show is that neither order of experience helps here. Adding the same information term to the
sender's objective gives an informative code in seven of the eight evaluated checkpoints: they reach 6.1 to 7.8\% of
new triples against controls of 0.7
to 1.4\%, with 29 to 42 names for 300 rules, modal-message purity 0.55 to 0.62, and the rule decoded at 2.3 to
3.8\%, above both chance (0.3\%) and the permutation null (0.7\%). The eighth collapsed like a joint run (1.7\%
against 1.9\%); over the eight, $6.3\pm0.7$\% against $1.3\pm0.1$\%. Where the code emerges it is informative but
much coarser than in the small networks, where reward-trained codes are decoded at 7 to 35\%.

For reference, an executor of the same model and adapter recipe was trained in a companion study to read the true
rule written in English (``swap the colours 9 and 7, then move everything right'') instead of a sender's symbols.
Evaluated with the procedure of this pilot, on the same programs and the same query episodes, it reaches 56 to
62\% exact match on the 91 unseen pairs and 40 to 43\% on the 231 new triples (three seeds), where the seven
informative emergent codes reach 8.4 to 9.9\% and 6.1 to 7.8\%.
This says that the model can execute a rule when given an informative description of it. It does not locate the
emergent system's remaining errors: that executor is given the rule that a sender has to infer, reads pretrained
English rather than an invented code, and was trained in a different configuration.

The pilot is small, and we report it as such. No statistical test is performed on these runs, and the
constant-message control, not the other-episode control, is the reference here: episodes are grouped by program
inside a batch, so a rotated message often belongs to the same rule. What the pilot shows is that the degenerate
solution and the remedy for it appear in a pretrained model as well; it does not show that the effect sizes match.

\begin{table}[h]\centering\small\resizebox{\textwidth}{!}{\begin{tabular}{lccccccccc}\toprule
 & \multicolumn{4}{c}{runs} & \multicolumn{2}{c}{new triples (\%)} & \multicolumn{3}{c}{the code} \\ \cmidrule(lr){2-5}\cmidrule(lr){6-7}\cmidrule(lr){8-10}
condition & started & finished & non-finite & averaged & own message & constant message & names & purity & rule decoded (\%) \\ \midrule
joint & 4 & 4 & 0 & 4 & 1.8 $\pm$ 0.2 & 1.8 $\pm$ 0.2 & 3 $\pm$ 2 & 0.90 $\pm$ 0.10 & 0.1 $\pm$ 0.0 \\
staged & 4 & 4 & 0 & 4 & 1.6 $\pm$ 0.2 & 1.6 $\pm$ 0.2 & 2 $\pm$ 0 & 0.98 $\pm$ 0.02 & 0.2 $\pm$ 0.1 \\
joint + information term & 8 & 8 & 4 & 8 & 6.3 $\pm$ 0.7 & 1.3 $\pm$ 0.1 & 36 $\pm$ 2 & 0.56 $\pm$ 0.03 & 2.7 $\pm$ 0.4 \\
\bottomrule
\end{tabular}
}
\caption{The same game played by a pretrained model (Qwen2.5-0.5B with low-rank adapters, 5\,000 steps). Runs are
counted at each stage; a run whose 16-bit adapters became non-finite is evaluated at its last snapshot before the
event, and averages use all eight. Accuracy is exact match on the 231
held-out triples no training program computes, beside the same executor's accuracy when every message is replaced
by a constant. Names are the distinct most frequent messages over the 300 held-out triples, purity the frequency of
a rule's most frequent message, and the last column how often the rule is decoded from the message on held-out
episodes (nearest neighbour; chance 0.3\%, permutation null 0.7\%). Mean $\pm$ s.e.m. over runs.}\label{tab:llm}
\end{table}


\section{Further detail on the collapse results}\label{s:detail}
These paragraphs carry the detail behind four statements of the main text: the controls that add variety without
information, the six other libraries, the order in which names are lost, and the human iterated-learning chains.

\paragraph{Controls that add variety without information.} Counting distinct messages overstates how differentiated
a code is, because a sender can vary its message between episodes while giving many rules the same name. We
therefore measured, for every run, how often the message identifies its rule on held-out episodes, by
nearest-neighbour decoding (Section~\ref{s:decoding}). Controls that make a code more
varied without making it much more informative restore little. A reward that pays only for the cells the rule
changes, and nothing for copying the input, raises the number of names for the thirty single operations from 6 or
7 to 13 to 16, but half the operations still share one name. Its messages identify a new triple in 10 to 11\% of
episodes, closer to the collapsed code's 8 to 10\% than to the staged code's 24 to 35\%, against 1\% by chance, and
it solves 1.9\% of new triples. The same reward brings the staged arm down from 15.8\% to 3.2\% (two runs).
Keeping the sender exploring for longer, by annealing its entropy floor from 2.5 to 0.3~nat over the first eight
thousand steps, adds a little of both, 16 to 18\% and 4.2 to 4.5\%.

\paragraph{Other libraries.} Whether joint experience collapses depends on the degenerate solution being attractive,
and we could not predict what makes it attractive. We repeated the manipulation in six more libraries
(Fig.~4e of the main text). Four are subsets of the original library: its first twenty-four operations, its first
seventeen, a random seventeen mixing all its families, and the library without the seven operations its frozen
perception confuses. The other two are a second library of seventeen operations sharing none with the first, and
that second library with six near-duplicates added. Every subset of the first library collapses under joint experience except
the random seventeen. The second library never collapses, with or without near-duplicates, and there staged
experience, joint experience and the added term land within a point of one another, 6.7 to 7.3\% on its new
triples, close to the real-valued channel for that library, 8.2\%. In two libraries the staged arm
varies widely between its two runs, 4.1 and 11.6\% with the first seventeen operations and 5.4 and 11.1\% with
the first twenty-four, against 2.3 and 5.7\% for joint experience and 12.7 and 17.2\% for the real-valued
channel. The staged advantage there is positive but not resolved with two runs per arm (5.5 points, $-0.1$ to
12.7, and 2.6 points, $-1.4$ to 8.0); in both libraries the more accurate staged run is also the one whose
messages identify the triple more often. Each comparison uses that library's own new triples, 49 to 178 of 300.

\paragraph{The order of collapse.} Collapse is ordered rather than uniform. Where it is partial, under staged experience or under joint
experience with the added term, the operations that lose a name of their own are first the ones the perception
confuses, the rotations and the transpose keeping a distinct name in 0 to 25\% of runs against 67 to 72\% for
clearly separable operations, and then the colour-changing ones. Where collapse is total, only the three size
operations keep a name. Perceptual confusability shapes which names go first, without deciding whether the code
collapses.

\paragraph{Human iterated learning.} In the iterated-learning experiment of Kirby, Cornish and Smith
\cite{kirby2008}, a participant learns an artificial language for twenty-seven meanings and their output becomes
the next participant's input. We re-analysed the published chains with the statistic of
Fig.~4c of the main text. Under plain transmission the language collapses from twenty-seven distinct signals to
between two and five within ten generations, 0.14 per meaning, one chain ending on a single word for twenty-five
meanings. When the experimenters filtered homonyms out of each training set, a pressure for expressivity, the
languages kept 0.59 signals per meaning and became compositional (difference at the tenth generation 0.45, 95\%
confidence interval 0.17 to 0.74; exact permutation $p = 0.029$, the smallest attainable with four chains per
condition; Fig.~5 of the main text). The human collapse is ordered too: without the pressure the languages lose
colour first, the information a signal carries about it falling to 0.06 by the tenth generation, then shape at
0.18, keeping motion at 0.66; with the pressure all three survive, at 0.63, 0.59 and 0.82. The pressures differ, a
filter on the training set there and an information term or staged experience here, and so do the mechanisms, a
learning bottleneck against exploration under reinforcement. What is shared is the abstract fact: a code shaped
only by the pull towards compression degenerates, and a pressure for expressivity keeps it differentiated, in
human learners and in our sender, measured the same way.

\end{document}
